\section{Conclusions, kill criteria, and the next decision}
\label{sec:tpc68-claim-boundary}

\subsection{What is proved}

For the complete literal missing-native map at one fixed nonzero shift:
\begin{enumerate}[label=\textup{(\arabic*)},leftmargin=3em]
 \item physical atomic normalization gives a positive Gram \(K=I+H\)
       with a Hermitian zero-diagonal alias defect;
 \item the common source/residual row gauge exposes the exact signs
       \(\mu(mj+h_0)\mu(nj+h_0)\) without changing any spectral quantity;
 \item on a same-residual row-pair collision this target sign freezes to
       \(\mu(d_m)\mu(d_n)\), so there is no pair-internal M\"obius
       oscillation in \(H_{\rm res}\); correspondingly, every pure
       \(H_{\rm res}\) closed walk has target-sign product one;
 \item sign-reversing partial matchings and ordered Abel sums give exact
       complex-weight pair certificates;
 \item row stars satisfy the exact sandwich
       \(s_*\le\|H\|\le\sqrt{\Delta_H}s_*\), and the local extremizer
       forbids replacing the star by a distinguished scalar row sum;
 \item a symmetric pair majorant gives the reducible-safe Perron optimum
       \(\|H\|\le\rho(P)\);
 \item even power rows and even traces retain cancellation at the level of
       complete decorated operator walks, but genuinely new signed
       information comes only from walks containing \(H_{\rm erase}\)
       edges or from literal complex-weight phases; and
 \item every certificate bounds both extremal eigenvalues of \(K\), while
       the full represented-plus-missing problem is exactly a direct Gram
       minus the positive nuisance penalty \(C_N\).
\end{enumerate}
These are L0 theorems.  Their identification with the actual fixed-\(h_0\)
missing carrier, complete native output labels, and raw atomic masses is
L1.

\subsection{What remains unproved}

The paper does not prove:
\begin{itemize}
 \item a power saving for one literal pair defect uniformly over all
       colliding active row pairs;
 \item an \(\ell^2\)-star saving that beats half the actual degree exponent;
 \item a fixed- or growing-length signed operator-walk estimate with
       constants uniform in every physical block;
 \item the required missing-mass saving or represented lower-frame
       comparison in the same scope;
 \item a bound \(C_N=o(1)\) for the nuisance shorting penalty; or
 \item any parity estimate, prime-pair theorem, or twin-prime conclusion.
\end{itemize}

\subsection{Kill criteria}

The following stops have different strengths.
\begin{enumerate}[label=\textup{(K\arabic*)},leftmargin=3.5em]
 \item \textbf{Wrong object.}  If a proposed proof erases
       \(\gamma\), side, external block, opposite row, or any label still
       orthogonal in the native range, or replaces complete outputs by
       reduced cells, it does not estimate the operator in this paper.
 \item \textbf{Wrong quantifier.}  If the carrier, partition, pairing, or
       row deletion depends on the later coefficient vector or extremizing
       eigenvector, it is not a coefficient-uniform certificate.
 \item \textbf{Wrong physical regime.}  A fixed-pair result, average over
       pairs, average over shifts, density-one row result, or frozen/smoothed
       family does not imply the required growing-row theorem at the
       prescribed fixed \(h_0\).
 \item \textbf{Pair-route stop.}  Failure of
       \(\tau\ge\delta\), or a pair majorant with large \(\rho(P)\), kills
       only the pair/Perron certificate.  It does not rule out cancellation
       between complete stars or operator walks.
 \item \textbf{Star-route stop.}  A star bound that does not beat
       the target after multiplication by \(\sqrt{\Delta_H}\)---equivalently,
       \(\sigma<\delta/2\) in the exponent model of
       \cref{prop:tpc68-conditional-star}---kills only the star
       certificate.  It does not rule out a later even-walk or trace
       estimate.
 \item \textbf{Walk-route stop.}  An unsigned support-walk count or a
       nonuniform fixed-moment estimate is not evidence for
       \eqref{eq:tpc68-L2-walk-hypothesis}.  Failure of one moment length
       does not kill every other length.
 \item \textbf{Hard uniform stop.}  If the actual shorted Gram
       \(G_{\rm sh}\) has a kernel, equivalently if some \(z\) satisfies
       \(Az\ne0\) and \((A+M)z=0\), then
       \(\delta(A,M)=0\).  No canonical-representative estimate can repair
       that uniform reassembly claim.  A distinguished physical vector
       remains a logically separate question.
 \item \textbf{Endpoint stop.}  If the total literal loss is at least
       \(1/400\), the current endpoint implication must be abandoned even
       if every intermediate qualitative statement is correct.
\end{enumerate}

\subsection{Route decision after TPC--68}

TPC--67 and the present paper now expose two independent finite-dimensional
doors on actual coefficients:
\[
 \text{all-row determinant mass}
 \qquad\text{and}\qquad
 \text{signed native operator walks}.
\]
The determinant door needs physical matching, pivot, and cycle-activity
input.  The signed-operator door needs one of the literal pair, star, or
walk premises above, followed by missing-mass and nuisance control.  A
useful next paper should attack one \emph{specified} fixed-\(h_0\) decorated
walk or star family and either prove a uniform saving or produce an actual
coefficient counterexample that closes that door.  More support-only
repackaging would not advance the proof chain.

\subsection{Status line}

TPC--68 closes a necessary operator interface and identifies the exact
location of the genuine two-affine signs after residual gauge removal.  It
does not break the arithmetic parity barrier.  A negative result for one
certificate family would still be publishable as a sharp route
obstruction, but it must not be presented as evidence for or against the
twin-prime conjecture itself.
